\section{Вычислительный эксперимент. Принципы проведения вычислительного эксперимента. Модель, алгоритм, программа}

\subsection{Понятие вычислительного эксперимента}

\textbf{Вычислительный эксперимент} --- исследование математической модели объекта или процесса с помощью вычислительных методов и компьютерных программ.

В отличие от натурного эксперимента, где исследуется непосредственно реальный объект, в вычислительном эксперименте исследуется его математическая модель.

В общем виде математическую модель можно записать как
\[
\mathcal{M}(u,p,x,t)=0,
\]
где
\begin{itemize}
	\item $u$ --- искомые характеристики системы;
	\item $p$ --- параметры модели;
	\item $x,t$ --- пространственные и временные переменные;
	\item $\mathcal{M}$ --- совокупность уравнений, ограничений, начальных и граничных условий.
\end{itemize}

После выбора численного метода непрерывная математическая задача заменяется конечномерной вычислительной задачей
\[
\mathcal{M}_h(u_h,p)=0,
\]
где $h$ обозначает параметры дискретизации: шаг сетки, шаг по времени, число элементов и т.п.

Таким образом, вычислительный эксперимент представляет собой последовательность вычислений, выполняемых для исследования поведения модели при различных исходных данных, параметрах и режимах.

Он особенно важен в тех случаях, когда натурный эксперимент:
\begin{itemize}
	\item слишком дорог;
	\item технически сложен;
	\item опасен;
	\item занимает слишком много времени;
	\item принципиально невозможен;
	\item не позволяет непосредственно наблюдать интересующие характеристики.
\end{itemize}


\subsection{Цели и задачи вычислительного эксперимента}

Основной целью вычислительного эксперимента является получение информации о свойствах исследуемой системы на основе её математической модели.

К основным задачам относятся:

\begin{enumerate}
	\item \textbf{Исследование поведения системы.}
	
	Определяется изменение характеристик объекта во времени и пространстве, устанавливаются характерные режимы его работы.
	
	\item \textbf{Исследование влияния параметров.}
	
	Определяется зависимость результата
	\[
	y = F(p_1,p_2,\ldots,p_m)
	\]
	от параметров модели.
	
	\item \textbf{Поиск критических режимов.}
	
	Исследуются условия потери устойчивости, возникновения резонанса, разрушения, перехода к другому режиму и т.п.
	
	\item \textbf{Оптимизация.}
	
	Требуется найти параметры $p$, минимизирующие или максимизирующие некоторый критерий
	\[
	J(p)\rightarrow \min_{p\in P}
	\]
	при выполнении ограничений модели.
	
	\item \textbf{Проверка гипотез.}
	
	С помощью серии расчётов сравниваются различные предположения о механизмах исследуемого процесса.
	
	\item \textbf{Прогнозирование.}
	
	Модель используется для предсказания поведения системы в условиях, которые ещё не наблюдались экспериментально.
	
	\item \textbf{Снижение объёма натурных экспериментов.}
	
	Вычисления позволяют предварительно определить наиболее интересные области параметров и сократить число дорогостоящих физических экспериментов.
\end{enumerate}

Важно, что вычислительный эксперимент не устраняет необходимость натурного эксперимента. Он дополняет его и позволяет исследовать модель в значительно более широком диапазоне условий.


\subsection{Триада «модель--алгоритм--программа»}

Основой вычислительного эксперимента является триада

\[
\boxed{\text{математическая модель}
	\longrightarrow
	\text{вычислительный алгоритм}
	\longrightarrow
	\text{программа}}
\]

Каждый из этих уровней отвечает на свой вопрос.

\subsubsection*{Математическая модель}

\textbf{Модель} описывает существенные свойства исследуемого объекта средствами математики.

Она может включать:
\begin{itemize}
	\item алгебраические уравнения;
	\item обыкновенные дифференциальные уравнения;
	\item уравнения в частных производных;
	\item интегральные уравнения;
	\item вероятностные зависимости;
	\item ограничения;
	\item начальные и граничные условия.
\end{itemize}

Например, процесс теплопроводности может описываться уравнением
\[
\frac{\partial u}{\partial t}
=
a^2\frac{\partial^2u}{\partial x^2}.
\]

Модель представляет физическую, экономическую, биологическую или иную предметную задачу в математической форме.

\subsubsection*{Алгоритм}

Как правило, математическая модель не может быть непосредственно решена компьютером.

Поэтому строится \textbf{вычислительный алгоритм} --- конечная последовательность операций, позволяющая получить приближённое решение задачи.

Например, производную можно заменить конечной разностью
\[
\frac{\partial^2u}{\partial x^2}(x_i,t_k)
\approx
\frac{u_{i-1}^{k}-2u_i^{k}+u_{i+1}^{k}}{h^2}.
\]

После дискретизации исходная задача сводится к конечномерной системе алгебраических уравнений.

Для вычислительного алгоритма существенны:
\begin{itemize}
	\item корректность;
	\item точность;
	\item устойчивость;
	\item сходимость;
	\item вычислительная сложность;
	\item требования к памяти;
	\item возможность распараллеливания.
\end{itemize}

\subsubsection*{Программа}

\textbf{Программа} представляет собой реализацию алгоритма на конкретной вычислительной системе.

Она включает:
\begin{itemize}
	\item представление исходных данных;
	\item реализацию вычислительных методов;
	\item организацию вычислений;
	\item хранение промежуточных результатов;
	\item ввод и вывод данных;
	\item средства анализа и визуализации результатов.
\end{itemize}

Важно различать ошибки трёх уровней.

Ошибочная модель не станет правильной при использовании точного алгоритма.

Правильная модель не даст правильного результата при некорректном численном методе.

Наконец, правильные модель и алгоритм могут быть испорчены ошибкой программной реализации.

Поэтому каждый уровень триады должен проверяться отдельно.


\subsection{Этапы проведения вычислительного эксперимента}

Типичный вычислительный эксперимент состоит из следующих этапов.

\begin{enumerate}
	\item \textbf{Постановка задачи.}
	
	Определяются объект исследования, цель моделирования, интересующие характеристики и диапазоны изменения параметров.
	
	\item \textbf{Построение математической модели.}
	
	Выбираются существенные свойства объекта, вводятся предположения и упрощения, формулируются уравнения, начальные и граничные условия.
	
	\item \textbf{Исследование математической задачи.}
	
	Желательно установить существование и единственность решения, свойства решения, характерные масштабы и возможные особые режимы.
	
	\item \textbf{Выбор численного метода.}
	
	Непрерывная задача заменяется дискретной. Выбираются сетка, шаги, аппроксимации и методы решения возникающих алгебраических систем.
	
	\item \textbf{Разработка вычислительного алгоритма.}
	
	Определяется последовательность операций, критерии завершения итераций, контроль ошибок и структура данных.
	
	\item \textbf{Программная реализация.}
	
	Алгоритм реализуется в виде программы или с использованием существующего программного комплекса.
	
	\item \textbf{Верификация.}
	
	Проверяется правильность вычислительного алгоритма и его программной реализации.
	
	\item \textbf{Проведение серии расчётов.}
	
	Выполняются вычисления для различных параметров, начальных данных и вариантов модели.
	
	\item \textbf{Валидация.}
	
	Результаты модели сопоставляются с экспериментальными данными или другими достоверными источниками.
	
	\item \textbf{Анализ результатов.}
	
	Определяются закономерности, чувствительность решения к параметрам, области устойчивости и границы применимости модели.
	
	\item \textbf{Уточнение модели.}
	
	При необходимости изменяются модель, численный метод или параметры, после чего вычислительный эксперимент повторяется.
\end{enumerate}

Таким образом, вычислительный эксперимент имеет итерационный характер:
\[
\text{модель}
\rightarrow
\text{расчёт}
\rightarrow
\text{сравнение}
\rightarrow
\text{уточнение модели}.
\]


\subsection{Планирование серии расчётов}

Отдельный расчёт обычно даёт мало информации. Для исследования модели проводится \textbf{серия вычислительных экспериментов}.

Пусть модель зависит от параметров
\[
p=(p_1,\ldots,p_m).
\]

Необходимо выбрать набор точек
\[
p^{(1)},p^{(2)},\ldots,p^{(N)},
\]
в которых будут выполнены расчёты.

При планировании определяются:
\begin{itemize}
	\item диапазоны изменения параметров;
	\item значения управляющих факторов;
	\item шаг изменения параметров;
	\item исследуемые выходные характеристики;
	\item критерии сравнения различных режимов;
	\item требуемая точность вычислений.
\end{itemize}

Наиболее простой подход --- изменение одного параметра при фиксированных остальных:
\[
y=y(p_i).
\]

Однако при большом количестве параметров такой подход становится неэффективным, поскольку число возможных комбинаций быстро возрастает.

Поэтому могут применяться методы \textbf{планирования эксперимента}, позволяющие получить максимальное количество информации при ограниченном числе расчётов.

Важной задачей является также \textbf{анализ чувствительности}. Для параметра $p_i$ можно рассматривать величину
\[
S_i =
\frac{\partial y}{\partial p_i},
\]
характеризующую чувствительность результата $y$ к изменению данного параметра.

Если вычисление производной затруднено, используется конечная разность
\[
S_i
\approx
\frac{y(p_i+\Delta p_i)-y(p_i)}
{\Delta p_i}.
\]

Планирование расчётов должно учитывать не только свойства модели, но и вычислительную стоимость одного запуска программы.


\subsection{Верификация алгоритма и программы}

\textbf{Верификация} отвечает на вопрос:

\begin{center}
	\textit{«Правильно ли мы решаем математическую задачу?»}
\end{center}

При верификации проверяется соответствие вычислительного результата исходной математической модели.

Следует различать несколько источников ошибок:
\[
\varepsilon_{\text{общ}}
\approx
\varepsilon_{\text{дискр}}
+
\varepsilon_{\text{итер}}
+
\varepsilon_{\text{окр}},
\]
где
\begin{itemize}
	\item $\varepsilon_{\text{дискр}}$ --- ошибка дискретизации;
	\item $\varepsilon_{\text{итер}}$ --- ошибка незавершённого итерационного процесса;
	\item $\varepsilon_{\text{окр}}$ --- ошибки округления.
\end{itemize}

Основные способы верификации:

\begin{enumerate}
	\item \textbf{Сравнение с аналитическим решением.}
	
	Если для специального случая известно точное решение $u$, вычисляется ошибка
	\[
	e_h=\|u-u_h\|.
	\]
	
	\item \textbf{Исследование сходимости по сетке.}
	
	Расчёт повторяется на последовательности сеток
	\[
	h,\quad \frac{h}{2},\quad \frac{h}{4}.
	\]
	
	Если метод имеет порядок $p$, то ожидается
	\[
	\|u-u_h\| = O(h^p).
	\]
	
	При отсутствии точного решения сравнивают результаты на различных сетках:
	\[
	\|u_h-u_{h/2}\|.
	\]
	
	\item \textbf{Проверка законов сохранения.}
	
	Например, если исходная модель сохраняет массу или энергию, необходимо контролировать соответствующий баланс.
	
	\item \textbf{Сравнение с независимой программой или методом.}
	
	Одна задача решается различными численными методами или программными пакетами.
	
	\item \textbf{Тестирование программы.}
	
	Проверяются отдельные программные модули, граничные случаи, обработка входных данных и устойчивость программы к некорректным данным.
\end{enumerate}

Таким образом, верификация относится прежде всего к связке
\[
\text{математическая модель}
\longrightarrow
\text{алгоритм}
\longrightarrow
\text{программа}.
\]


\subsection{Валидация математической модели}

\textbf{Валидация} отвечает на другой вопрос:

\begin{center}
	\textit{«Правильно ли математическая модель описывает реальный объект?»}
\end{center}

Для этого результаты вычислений сопоставляются с экспериментальными данными
\[
y_i^{\text{calc}}
\quad\text{и}\quad
y_i^{\text{exp}}.
\]

Например, можно использовать среднеквадратичное отклонение
\[
E =
\sqrt{
	\frac{1}{N}
	\sum_{i=1}^{N}
	\left(
	y_i^{\text{calc}}-
	y_i^{\text{exp}}
	\right)^2
}.
\]

Валидация может включать:
\begin{itemize}
	\item качественное сравнение характера процессов;
	\item количественное сравнение с экспериментом;
	\item сравнение характерных величин;
	\item проверку предсказаний модели;
	\item исследование диапазона её применимости.
\end{itemize}

Необходимо чётко различать верификацию и валидацию:

\[
\boxed{
	\begin{array}{ll}
		\text{верификация:} & \text{правильно ли решены уравнения?}\\[2mm]
		\text{валидация:} & \text{правильные ли уравнения выбраны для объекта?}
\end{array}}
\]

Даже полностью верифицированная программа может давать физически неправильные результаты, если сама математическая модель неадекватна объекту.


\subsection{Анализ и представление результатов}

После выполнения расчётов необходимо не просто получить набор чисел, а извлечь из него информацию о свойствах исследуемой системы.

Анализ может включать:
\begin{itemize}
	\item построение зависимостей выходных величин от параметров;
	\item исследование пространственных и временных распределений;
	\item определение максимумов и минимумов;
	\item выявление критических значений параметров;
	\item исследование устойчивости;
	\item анализ чувствительности;
	\item статистическую обработку результатов;
	\item сравнение нескольких моделей или алгоритмов.
\end{itemize}

Результаты удобно представлять в виде:
\begin{itemize}
	\item таблиц;
	\item графиков;
	\item полей и изолиний;
	\item трёхмерных визуализаций;
	\item диаграмм;
	\item анимаций динамических процессов.
\end{itemize}

При представлении результатов обязательно должны быть указаны:
\begin{itemize}
	\item значения параметров модели;
	\item используемый численный метод;
	\item параметры дискретизации;
	\item критерии остановки итераций;
	\item единицы измерения;
	\item оценка точности результата.
\end{itemize}

В противном случае результат трудно интерпретировать и практически невозможно воспроизвести.


\subsection{Воспроизводимость вычислительного эксперимента}

\textbf{Воспроизводимость} означает возможность повторить вычислительный эксперимент и получить те же или статистически согласующиеся результаты.

Для воспроизводимости необходимо сохранять:

\begin{itemize}
	\item точную формулировку математической модели;
	\item исходные данные;
	\item значения параметров;
	\item начальные и граничные условия;
	\item используемый численный метод;
	\item параметры дискретизации;
	\item версию программы;
	\item версии используемых библиотек;
	\item параметры вычислительной среды;
	\item случайные начальные значения при использовании стохастических методов.
\end{itemize}

Для стохастического алгоритма, например,
\[
X^{(k+1)}=F(X^{(k)},\xi_k),
\]
результат зависит от последовательности случайных чисел $\xi_k$. Поэтому для точного повторения расчёта необходимо фиксировать начальное состояние генератора случайных чисел.

В современных вычислительных исследованиях желательно также использовать:
\begin{itemize}
	\item системы контроля версий;
	\item автоматизированные сценарии запуска расчётов;
	\item конфигурационные файлы вместо ручного задания параметров;
	\item автоматическую обработку результатов;
	\item сохранение метаданных вычислительного эксперимента.
\end{itemize}

Воспроизводимость позволяет отличить научный вычислительный эксперимент от разового запуска программы, результат которого невозможно независимо проверить.

В итоге вычислительный эксперимент представляет собой не просто численное решение одной задачи, а полный цикл исследования:
\[
\boxed{
	\text{постановка задачи}
	\rightarrow
	\text{модель}
	\rightarrow
	\text{алгоритм}
	\rightarrow
	\text{программа}
	\rightarrow
	\text{расчёт}
	\rightarrow
	\text{верификация и валидация}
	\rightarrow
	\text{анализ результатов}
}
\]

\newpage